\section{数学是被发现的，还是被创造的}

访谈早段讨论“被创造的与被发现的”。这是 AI for Math 的哲学底色。若数学是被发现的，AI 需要学习客观结构：素数、对称、拓扑、曲率和证明之间的稳定关系。若数学是被创造的，AI 需要学习人类如何选择定义、公理、记号和理论语言。现实中两者交织：数学对象有客观约束，数学语言和证明路径又由人创造。

\begin{figure}[H]
\centering
\includegraphics[width=0.92\textwidth]{math-created-discovered.png}
\caption{数学的两种视角：被发现的结构与被创造的语言。}
\end{figure}

\begin{knowledgebox}{读图：为什么 AI for Math 同时面对两层任务}
左侧强调数学结构的客观性，右侧强调数学语言和公理系统的创造性。AI for Math 不能只做模式匹配，也不能只生成漂亮文字；它既要捕捉稳定结构，又要把推理写成可验证语言。
\end{knowledgebox}

\subsection{bounded attention 与 free attention}

访谈中出现 bounded vs free attention 的说法，可以理解为数学思考的两种模式：bounded attention 围绕定义、条件、目标和已有工具局部推进；free attention 则允许联想、审美、类比和跳跃。数学研究需要在二者之间切换：形式化证明偏 bounded，猜想和直觉生成偏 free。

\begin{figure}[H]
\centering
\includegraphics[width=0.92\textwidth]{bounded-free-attention.png}
\caption{数学直觉中的 bounded attention 与 free attention。}
\end{figure}

\begin{knowledgebox}{读图：两种注意力分别服务什么}
Bounded attention 适合证明检查、局部推理和 Lean 形式化；free attention 适合发现模式、提出猜想和跨领域类比。一个强 AI for Math 系统需要二者协同，而不是只会严格搜索或只会发散联想。
\end{knowledgebox}

\begin{warningbox}{不要把直觉和形式化对立起来}
数学直觉不是形式化的敌人。好的形式化系统把直觉变成可检查对象；好的直觉则帮助选择值得形式化的问题和证明路线。
\end{warningbox}

\subsection{本章小结}

AI for Math 的难点来自双重结构：数学既像客观世界，又像人类语言工程。模型必须在自由联想和严格验证之间来回切换。

